\section{Multimodal Foundation Agents: VIMA, RT-2 y RoboCat}

La ruta del video amplía los datos; la otra ruta unifica la interfaz. \term{multimodal prompting (prompting multimodal)} se refiere a que en el prompt se intercalan texto, imágenes, fotogramas objetivo, demostraciones en video u otros tokens perceptuales, y el modelo produce una action a partir de ellos. No consiste en añadir una imagen decorativa a un prompt de lenguaje, sino en hacer que los objetos visuales, las restricciones y las demostraciones formen parte directa de la especificación de la tarea.

\subsection{Unificar Text, Image, Video, Audio y Action}

Un multimodal foundation model intenta proyectar entradas distintas a un espacio de secuencias compartido y luego generar lenguaje o action según el contexto. Para un agent, la ventaja es que la tarea puede expresarse como «pon este objeto en la posición que muestra esta imagen» o mediante «una demostración»; el reto es que cada modality tiene requisitos distintos de tokenización, temporalidad, precisión y safety.

\lecturefigure{slide-50-multimodal-foundation-model.jpg}{Visión de un foundation model que unifica texto, imagen, video, sonido y action}{Diapositiva V050 recuperada de la grabación oficial; 00:43:03}

\begin{knowledgebox}{Lectura de la figura: un token space unificado no equivale a un significado físico unificado}
En la figura, todas las modality entran en un solo modelo, pero el costo de un error es distinto para un token de texto, un patch visual y una robot action. Primero observe cómo se codifica la entrada y después si la salida está sujeta a una action-space constraint; una architecture unificada favorece la knowledge transfer, pero sigue necesitando encoder específicos de cada modality, controller y safety validation.
\end{knowledgebox}

\lecturefigure{slide-51-multimodal-prompting-robots.jpg}{Una tarea robótica puede definirse con un prompt que intercala elementos visuales y de lenguaje}{Diapositiva V051 recuperada de la grabación oficial; 00:43:48}

El prompt puede incluir imágenes del objeto objetivo, esquemas de posición, zonas prohibidas y relaciones en lenguaje natural. Al leer la figura, primero distinga «a qué objeto se refiere», «cuál es el estado objetivo» y «qué constraint debe cumplir la acción», y después observe la robot trajectory. El prompt multimodal aumenta la specification bandwidth, pero también introduce ambiguity; el modelo debe hacer grounding en la scene actual, en lugar de emparejar plantillas de apariencia vistas en el entrenamiento.

\subsection{VIMA: Visual Manipulation with Multimodal Prompts}

VIMA introduce en un Transformer tokens de objetos visuales intercalados con tokens de texto y predice la robot action de forma autoregressive. Su aporte principal no es un manipulation benchmark aislado, sino expresar con un mismo model/interface familias de tareas distintas, como visual goal reaching, novel concept grounding, one-shot demonstration y constraint satisfaction.

\lecturefigure{slide-52-vima-architecture.jpg}{VIMA modela de forma unificada el multimodal prompt, la observation y la action}{Diapositiva V052 recuperada de la grabación oficial; 00:44:33}

\begin{knowledgebox}{Lectura de la figura: la conexión entre el prompt encoder y el action decoder}
Primero observe cómo las imágenes y el texto del prompt forman la token sequence; luego, cómo la observation actual se incorpora al contexto, y por último, cómo el action decoder produce tokens de pose o de control. El Transformer compartido de VIMA favorece la transferencia entre tareas, pero el tabletop world, el object set y la action representation del benchmark siguen limitando el alcance de su evidencia.
\end{knowledgebox}

Si los prompt tokens son $p_{1:m}$, las observation históricas son $o_{1:t}$ y la action sequence es $a_{1:t}$, el autoregressive objective puede escribirse como:
\[
\mathcal{L}_{\mathrm{VIMA}}(\theta)
=-\sum_{t=1}^{T}\log p_\theta\!\left(a_t\mid p_{1:m},o_{1:t},a_{<t}\right).
\]
Aquí $p_{1:m}$ puede intercalar texto y objetos visuales, $o_{1:t}$ son las observation actual e históricas, $a_{<t}$ son las acciones previas, $T$ es el episode horizon y $\theta$ son los parámetros del modelo. La fórmula expresa un conditional sequence modeling unificado, pero no implica que todo action error equivalga a un error en un token de lenguaje.

\subsection{Límites de la evidencia para las cuatro capacidades de Prompt}

Las cuatro diapositivas siguientes evalúan, respectivamente, objetivos visuales, conceptos temporales, una sola demostración y restricciones. Deben verse como formas distintas de task specification, no como cuatro demo sin relación entre sí; al leer las figuras, verifique continuamente qué información aporta el prompt, sobre qué objeto debe hacer grounding el modelo y si la evaluación se realiza solo dentro de una tabletop distribution controlada.

\lecturefigure{slide-53-vima-visual-goal-reaching.jpg}{VIMA: una imagen objetivo especifica el estado visual que se desea alcanzar}{Diapositiva V053 recuperada de la grabación oficial; 00:44:51}

El objetivo visual reduce la carga de describir relaciones espaciales con lenguaje. Primero compare la scene inicial con la image objetivo y luego verifique las posiciones relativas de los objetos después de la acción; el éxito indica que el modelo puede convertir la apariencia objetivo en una manipulation sequence, pero por sí solo no demuestra que comprenda la causalidad física en general ni las oclusiones del mundo real.

\lecturefigure{slide-54-vima-concept-grounding.jpg}{VIMA: grounding de conceptos nuevos definidos temporalmente dentro del prompt}{Diapositiva V054 recuperada de la grabación oficial; 00:45:03}

Palabras sin significado establecido, como «wug» y «blicket», se definen temporalmente mediante su correspondencia visual en el prompt. Al leer la figura, primero identifique el binding entre token y objeto, y luego vea si el modelo traslada la relación a la action; esto respalda el in-context concept grounding, no demuestra que el modelo haya visto el significado de esas palabras durante el preentrenamiento.

\lecturefigure{slide-55-vima-one-shot-demonstration.jpg}{VIMA: inferir la tarea a partir de una sola demostración visual y reproducirla}{Diapositiva V055 recuperada de la grabación oficial; 00:45:15}

Una sola demostración aporta una trajectory-level specification. Primero observe los cambios de los objetos en la demonstration y luego si la test scene tiene posiciones o instancias nuevas; el éxito exige que el modelo extraiga relaciones en lugar de copiar coordenadas absolutas. Que una sola demo pueda generalizarse sigue dependiendo de la similitud de la familia de tareas, la cámara y la object distribution.

\lecturefigure{slide-56-vima-constraint-satisfaction.jpg}{VIMA: completar una manipulation bajo restricciones visuales}{Diapositiva V056 recuperada de la grabación oficial; 00:45:48}

El constraint prompt le indica al agent qué zona, objeto o trayectoria no debe tocar. Al leer la figura, no basta con verificar el objetivo final: también hay que comprobar si la trajectory viola restricciones intermedias. Esto muestra que la agent evaluation debe incluir la seguridad del proceso, y no solo el final frame. Las restricciones en una mesa controlada siguen siendo distintas de la force, la collision y la human safety de un robot real.

\teachervoice{El ponente dice que VIMA «tokenizes everything» y que, con interleaved image--text tokens, unifica tareas que antes requerían un pipeline especializado. Esto debe entenderse como una unificación de la modeling interface, no como la desaparición del problema de control físico; el grounding visual, el action decoder y las restricciones del entorno siguen asumiendo cada uno su propia responsabilidad.}

\subsection{RT-2 y RoboCat: ampliar los Robot Data y la Transfer}

Al final, el curso usa RT-2 y RoboCat para mostrar ejemplos posteriores de la misma recipe. RT-2 hace co-fine-tune del web-scale vision-language knowledge junto con robot demonstration, y representa la action como tokens que pueden generarse; RoboCat entrena una policy unificada que abarca distintas tareas y robot embodiment. Ambos amplían la evidencia de transfer, pero siguen dependiendo de la mezcla de datos concreta y del action support.

\lecturefigure{slide-57-rt2-robocat.jpg}{RT-2 y RoboCat extienden la ruta del multimodal generalist robot}{Diapositiva V057 recuperada de la grabación oficial; 00:46:12}

\begin{knowledgebox}{Términos clave: tres interfaces de generalist-agent}
\begin{tabularx}{\textwidth}{@{}L{0.17\textwidth}L{0.25\textwidth}X@{}}
\toprule
Sistema & Interfaz principal & Mecanismo central y límites \\
\midrule
Voyager & JavaScript code-as-action & Horizon largo, componible; depende de una symbolic API \\
Eureka & LLM-generated reward code & Aprende control continuo con RL; depende del simulator y de la reward audit \\
VIMA / RT-2 / RoboCat & De multimodal token a robot action & Representación compartida y transferencia de datos; depende del action grounding y de la embodiment coverage \\
\bottomrule
\end{tabularx}
\end{knowledgebox}

\teachervoice{El ponente explica RT-2 como un modelo preentrenado primero con internet-scale data y luego sometido a fine-tune con human-collected robot demonstrations; RoboCat, por su parte, muestra una sola policy que abarca distintos robots y tareas. En clase se presentan como un promising follow-up, sin afirmar que la generalist robotics ya esté resuelta.}

\subsection{Resumen de la sección}

El multimodal prompting amplía la specification de la tarea: texto, imágenes objetivo, conceptos temporales, demostraciones y constraint pueden entrar en un contexto unificado. VIMA demuestra que varias clases de tabletop task pueden expresarse con un mismo Transformer/action interface, y RT-2 y RoboCat amplían además el web prior, los robot data y la embodiment transfer. Un modelo unificado aporta representaciones compartidas, pero el grounding, el control, la safety y el evidence boundary todavía deben verificarse por separado.
